% !TEX root = 00_MidReview.tex

\chapter{数学工具}

\section{线性差分系统}

\subsection{求解一阶线性差分方程}

考虑如下差分方程\sidenote{通解由齐次解和特解构成：求齐次解使用特征方程法；求特解用猜测形式法}
\begin{align*}
    y\nt+a\,y_t=C
\end{align*}
写出特征方程
\begin{align*}
    \lambda+a=0
\end{align*}
解得特征根 \(\lambda=-a\)，差分方程的齐次解 \(y^c_t=A\,\bka{-a}^t\)，其中 \(A\) 是任意常数。

猜测特解 \(y^p_t=k\) 代入差分方程
\begin{align*}
    k+a\,k=C \qquad\Rightarrow\qquad k=\frac{C}{1+a}
\end{align*}
结果要求 \(a\ne-1\)。单独考虑 \(a=-1\) 的情况，猜测特解 \(y^p_t=k\,t\) 代入差分方程
\begin{align*}
    k\,(t+1)-k\,t=C \qquad\Rightarrow\qquad k=C
\end{align*}
因此差分方程的通解为
\begin{equation*}
    y_t=y_t^c+y_t^p=\begin{cases}
        \ A\bka{-a}^{t}+\dfrac C{1+a} & a\neq -1\\
        \ A+C\,t & a=-1
    \end{cases}
\end{equation*}
假设系统的初始位置为 \(y_0\)，代入上式求解 \(A\) 的值
\begin{equation*}
    y_0=\begin{cases}
        \ A\bka{-a}^0 + \dfrac C{1+a} & a\neq -1\\
        \ A+C\,t & a=-1
    \end{cases} \quad \Longrightarrow \quad A=\begin{cases}
        \ y_0-\dfrac C{1+a} & a\neq -1\\
        \ y_0 & a=-1
    \end{cases}
\end{equation*}

\subsection{求解二阶线性差分方程}

考虑如下差分方程
\begin{align*}
    y\nnt+a\,y\nt+b\,y_t=C
\end{align*}
写出特征方程
\begin{align*}
    \lambda^2+a\,\lambda+b=0
\end{align*}
这是一元二次方程，判别式 \(\Delta=a^2-4\,b\) 的符号决定了特征根的类型，影响到差分方程齐次解的形式：
\begin{equation*}
    y^c=\begin{cases}
        \ A_1\,b^t_1+A_2\,b^t_2 & \n{ if}\quad \Delta>0\\
        \ A_3\,b^t + A_4\,t\,b^t & \n{ if}\quad \Delta=0\\
        \ \rho^t\bks{A_5\,\cos\bka{t\,\theta} + A_6\,\sin\bka{t\,\theta}} & \n{ if}\quad \Delta<0
    \end{cases}
\end{equation*}
再分别假设特解 \(y^p=k\)、\(y^p=k\,t\)、\(y^p_t=k\,t^2\) 代入差分方程解得 \footnote{2026 年 4 月 21 日勘误：此前第二种情况表达式有误}
\begin{equation*}
    y^p=\begin{cases}
       \ \dfrac C{1+a_1+a_2} & \n{ if}\quad 1+a_1+a_2\neq 0\\[3ex]
       \ \dfrac {C\,t}{2+a_1} & \n{ if}\quad 1+a_1+a_2=0\ \n{ and }\ a_1\neq-2\\[3ex]
       \ \dfrac {C\,t^2 }{2}& \n{ if}\quad a_1=-2\ \n{ and }\ a_2=1
    \end{cases}
\end{equation*}
根据方程的具体形式组合需要的齐次解和特解得到通解，匹配的条件已在公式中给出，不可随意组合以上结果。

\subsection{差分系统的稳定性}

\noxmubluebf{一阶系统的稳定性}

回顾一阶线性系统通解
\begin{equation*}
    y_t=\begin{cases}
        \ A\bka{-a}^{t}+\dfrac C{1+a} & a\neq -1\\
        \ A+C\,t & a=-1
    \end{cases}
\end{equation*}
根据指数函数 \(\bka{-a}^t\) 的收敛性，我们可以得出如下结论
\begin{itemize}
    \item 当 \(\bkf{a}<1\) 时，\(b^t\to 0\)，\(y_t\to C\,/\bka{1-a}\)，系统具有稳态 \(\bar y = C\,/\bka
    {1-a}\)\sidenote{\sidemath{\(\bkf{a}<1\)} 是差分系统具有稳态的充分条件，含 \sidemath{\(a=0\)}}
    \item 当 \(\bkf{a}>1\) 时，\(\bka{-a}^t\to \infty\)，\(y_t\to \infty\)，系统不具有稳态
    \item 当 \(a=-1\) 时，\(y_t=A+C\,t\)，系统不具有稳态
    \item 当 \(a=1\) 时，\(y_t=\bka{-1
    }^tA\)，系统不具有稳态
\end{itemize}
更严谨地说，以上 “不具有稳态” 要求系统的初始位置不在稳态上及其他值得忽略的情况。如果 \(a>0\)，\(y_t-\bar y\) 的符号会随着 \(t\) 交替变化，系统交替震荡收敛或交替震荡发散。

\myskip
\noxmubluebf{二阶系统的稳定性}

改写二阶线性差分方程，令 \(x_t=y\nt\)，有
\begin{equation*}
    \begin{cases}
        \ x\nt + a_1\,x_t + a_2\,y_t = C\\
        \ y\nt - x_t = 0
    \end{cases}
\end{equation*}
用矩阵符号表示为
\begin{align}
    \underbrace{\matbks{x\nt\\y\nt}}_{\bs X\nt}&=\underbrace{\matbks{-a_1 & -a_2\\1 & 0}}_{\bs A}\underbrace{\matbks{x_t\\y_t}}_{\bs X_t}+\underbrace{\matbks{C\\0}}_{\bs C}\\
    \bsX\nt&=\bsA\,\bsX_t+\bs C\label{eq:DiffEquMatrix}
\end{align}
关注矩阵 \(\bsA\)，设其迹为 \(T\)，行列式为 \(D\)，则特征方程为
\begin{align*}
    p\bka\lambda&=\bka{\lambda-\lambda_1}\bka{\lambda-\lambda_2}\\
    &= \lambda^2 - \bka{\lambda_1+\lambda_2}\lambda + \lambda_1\,\lambda_2\\
    &= \lambda^2 - T\,\lambda + D=0
\end{align*}
根据以下四个关键边界方程划分 \(\bka{T,D}\) 平面
\begin{align*}
    &(\n a)\qquad \Delta=T^2-4D=0\\
    &(\n b)\qquad D=\lambda_1\cdot\lambda_2=1\\
    &(\n c)\qquad p\bka1=1-T+D=0\\
    &(\n d)\qquad p\bka{-1}=1+T+D=0
\end{align*}
由 (b)、(c)、(d) 围出的 (6) 和 (7) 区域 —— \( \ p\bka{1},\,p\bka{-1}>0\) 且 \(D>1\) 为稳定区域，无论系统初始位置在何处，都会收敛至稳态。此外同样值得关注的是 (1) 和 (3) 区域 —— \( \ p\bka{1}\) 和 \(p\bka{-1}\) 异号，即 \(\bkf T>\bkf{1+D}\)，这样的系统具有鞍路径，存在两条（同一直线的两个方向）收敛路径。

\begin{figure}[H]
    \centering
    \sidenote{落在边界上也是不可能稳定的（排除特殊情况）；箭头表示大于的方向，如 \sidemath{\(p\bka{1}>1\)}}
    \begin{tikzpicture}[
        scale=1.0,
        line cap=round,
        line join=round,
        >=Latex
    ]
        % ---- User-tunable geometry ----
        \def\xmin{-5.2}
        \def\xmax{5.4}
        \def\ymin{-1.85}
        \def\ymax{4.1}
        \def\parabL{-4.00}
        \def\parabR{4.00}
        \def\lineL{-4.25}
        \def\lineR{4.25}
        \def\BaseLineWidth{0.85pt}
        \def\GuideLineWidth{0.60pt}
        
        \def\MirrorArrowLineWidth{1.00pt}
        \def\MirrorArrowALength{0.40}
        \def\MirrorArrowBLength{0.35}
        \def\MirrorArrowCLength{0.55}
        \def\MirrorArrowAY{3.35}
        \def\MirrorArrowBY{2.75}
        \def\MirrorArrowCY{1.00}
        \def\MirrorArrowAHalfGap{3.70}
        \def\MirrorArrowBHalfGap{3.75}
        \def\MirrorArrowCHalfGap{1.20}
        \def\MirrorArrowAAngle{-40}
        \def\MirrorArrowBAngle{135}
        \def\MirrorArrowCAngle{90}

        % ---- User-tunable number colors ----
        \colorlet{NumOneColor}{red!75!black}
        \colorlet{NumTwoColor}{red!75!black}
        \colorlet{NumThreeColor}{red!75!black}
        \colorlet{NumFourColor}{red!75!black}
        \colorlet{NumFiveColor}{red!75!black}
        \colorlet{NumSixColor}{red!75!black}
        \colorlet{NumSevenColor}{red!75!black}
        
        % ---- User-tunable label positions ----
        \def\pOneLabelX{-3.95}     \def\pOneLabelY{2.76}
        \def\pMinusOneLabelX{3.95} \def\pMinusOneLabelY{2.76}
        \def\DeltaLabelX{2.60}     \def\DeltaLabelY{3.85}
        \def\DOneLabelX{5.00}     \def\DOneLabelY{0.70}

        \def\labOneX{2.00}   \def\labOneY{-1.00}
        \def\labThreeX{-2.00}\def\labThreeY{-1.00}
        \def\labOneXX{3.00}   \def\labOneYY{0.50}
        \def\labThreeXX{-3.00}\def\labThreeYY{0.50}
        \def\labOneXXX{3.50}   \def\labOneYYY{1.70}
        \def\labThreeXXX{-3.50}\def\labThreeYYY{1.70}

        \def\labTwoX{0.35}   \def\labTwoY{-1.70}
        \def\labFiveX{1.50}  \def\labFiveY{2.72}
        \def\labSixX{0.85}   \def\labSixY{0.65}
        \def\labSevenX{0.35}\def\labSevenY{-0.27}

        \def\labFourX{4.20} \def\labFourY{3.60}

        \pgfmathsetmacro{\xcutL}{\ymin+1}
        \pgfmathsetmacro{\xcutR}{-1-\ymin}
        \pgfmathsetmacro{\MirrorArrowADx}{\MirrorArrowALength*cos(\MirrorArrowAAngle)}
        \pgfmathsetmacro{\MirrorArrowADy}{\MirrorArrowALength*sin(\MirrorArrowAAngle)}
        \pgfmathsetmacro{\MirrorArrowBDx}{\MirrorArrowBLength*cos(\MirrorArrowBAngle)}
        \pgfmathsetmacro{\MirrorArrowBDy}{\MirrorArrowBLength*sin(\MirrorArrowBAngle)}
        \pgfmathsetmacro{\MirrorArrowCDx}{\MirrorArrowCLength*cos(\MirrorArrowCAngle)}
        \pgfmathsetmacro{\MirrorArrowCDy}{\MirrorArrowCLength*sin(\MirrorArrowCAngle)}

        % Axes
        \draw[-{Latex[length=2.8mm,width=2mm]}, line width=\BaseLineWidth] (\xmin,0) -- (\xmax,0) node[below] {\(T\)};
        \draw[-{Latex[length=2.8mm,width=2mm]}, line width=\BaseLineWidth] (0,\ymin) -- (0,\ymax) node[right] {\(D\)};

        % Reference lines
        \draw[line width=\BaseLineWidth, black] (\lineL,{-\lineL-1}) -- (\xcutR,\ymin);
        \draw[line width=\BaseLineWidth, black] (\xcutL,\ymin) -- (\lineR,{\lineR-1});
        \draw[black, line width=\BaseLineWidth] (-5,1) -- (5,1);
        \draw[black, line width=\BaseLineWidth, domain=\parabL:\parabR, samples=320] plot (\x, {(\x*\x)/4});

        % Tick marks and labels
        \foreach \x/\lab in {-2/-2,-1/-1,1/1,2/2} {
            \draw (\x,0.06) -- (\x,-0.06);
            \node[below] at (\x,0) {\lab};
        }
        \draw (0.06,-1) -- (-0.06,-1);
        \node[right] at (0,-1) {\(-1\)};
        \draw (0.06,1) -- (-0.06,1);
        \node[above right] at (0,1) {\(1\)};

        % Helper guides at T = ±2
        \draw[black, densely dashed, line width=\GuideLineWidth] (-2,0) -- (-2,1);
        \draw[black, densely dashed, line width=\GuideLineWidth] (2,0) -- (2,1);

        % Labels for loci
        \node[black, anchor=west] at (\pMinusOneLabelX,\pMinusOneLabelY) {\(p\bka{-1}=0\)};
        \node[black, anchor=east] at (\pOneLabelX,\pOneLabelY) {\(p\bka{1}=0\)};
        \node[black] at (\DeltaLabelX,\DeltaLabelY) {\(\Delta=0\)};
        \node[black, left] at (\DOneLabelX,\DOneLabelY) {\(D=1\)};

        % Circled region numbers
        \newcommand{\RegionMarker}[4]{%
            \draw[color=#4, line width=0.55pt] (#1,#2) circle (0.15);
            \node[text=#4, font=\scriptsize\bfseries] at (#1,#2) {#3};
        }
        \RegionMarker{\labOneX}{\labOneY}{1}{NumOneColor}
        \RegionMarker{\labTwoX}{\labTwoY}{2}{NumTwoColor}
        \RegionMarker{-\labTwoX}{\labTwoY}{2}{NumTwoColor}
        \RegionMarker{\labThreeX}{\labThreeY}{3}{NumThreeColor}
        \RegionMarker{\labFourX}{\labFourY}{4}{NumFourColor}
        \RegionMarker{\labFiveX}{\labFiveY}{5}{NumFiveColor}
        \RegionMarker{\labSixX}{\labSixY}{6}{NumSixColor}
        \RegionMarker{\labSevenX}{\labSevenY}{7}{NumSevenColor}
        \RegionMarker{-\labFiveX}{\labFiveY}{5}{NumFiveColor}
        \RegionMarker{-\labSixX}{\labSixY}{6}{NumSixColor}
        \RegionMarker{-\labSevenX}{\labSevenY}{7}{NumSevenColor}
        \RegionMarker{-\labFourX}{\labFourY}{4}{NumFourColor}
        \RegionMarker{\labOneXX}{\labOneYY}{1}{NumOneColor}
        \RegionMarker{\labThreeXX}{\labThreeYY}{3}{NumThreeColor}
        \RegionMarker{\labOneXXX}{\labOneYYY}{1}{NumOneColor}
        \RegionMarker{\labThreeXXX}{\labThreeYYY}{3}{NumThreeColor}

        % User-tunable mirrored arrows
        \draw[-{Latex[length=2.4mm,width=1.7mm]}, black, line width=\MirrorArrowLineWidth]
            (\MirrorArrowAHalfGap,\MirrorArrowAY)
            -- ({\MirrorArrowAHalfGap+\MirrorArrowADx},{\MirrorArrowAY+\MirrorArrowADy});
        \draw[-{Latex[length=2.4mm,width=1.7mm]}, black, line width=\MirrorArrowLineWidth]
            (-\MirrorArrowAHalfGap,\MirrorArrowAY)
            -- ({-\MirrorArrowAHalfGap-\MirrorArrowADx},{\MirrorArrowAY+\MirrorArrowADy});
        \draw[-{Latex[length=2.4mm,width=1.7mm]}, black, line width=\MirrorArrowLineWidth]
            (\MirrorArrowBHalfGap,\MirrorArrowBY)
            -- ({\MirrorArrowBHalfGap+\MirrorArrowBDx},{\MirrorArrowBY+\MirrorArrowBDy});
        \draw[-{Latex[length=2.4mm,width=1.7mm]}, black, line width=\MirrorArrowLineWidth]
            (-\MirrorArrowBHalfGap,\MirrorArrowBY)
            -- ({-\MirrorArrowBHalfGap-\MirrorArrowBDx},{\MirrorArrowBY+\MirrorArrowBDy});
        \draw[-{Latex[length=2.4mm,width=1.7mm]}, black, line width=\MirrorArrowLineWidth]
            (\MirrorArrowCHalfGap,\MirrorArrowCY)
            -- ({\MirrorArrowCHalfGap+\MirrorArrowCDx},{\MirrorArrowCY+\MirrorArrowCDy});
        \draw[-{Latex[length=2.4mm,width=1.7mm]}, black, line width=\MirrorArrowLineWidth]
            (-\MirrorArrowCHalfGap,\MirrorArrowCY)
            -- ({-\MirrorArrowCHalfGap-\MirrorArrowCDx},{\MirrorArrowCY+\MirrorArrowCDy});
    \end{tikzpicture}
    \caption{\(\bka{T,D}\) 平面中的稳定性分类示意图}
    \label{fig:TDStabilityMap}
\end{figure}

\myskip
\noxmubluebf{Blanchard-Kahn 条件}

经济系统中存在两类变量：前定变量，在决策时已经确定的变量；跳跃变量，在决策时还未确定，可以自由选择的变量。BK 条件称，如果
\begin{itemize}
    \item \# 不稳定根 = \# 跳跃变量，存在唯一均衡路径（鞍路径）
    \item \# 不稳定根 > \# 跳跃变量，不存在稳定均衡路径
    \item \# 不稳定根 < \# 跳跃变量，稳定均衡路径不唯一
\end{itemize}



\newpage
\section{对数线性化}

\subsection{对数线性化的方法}

对数线性化常用本小节介绍的三种方法，各有优缺点，应在具体情况中灵活使用\sidenote{建议掌握多种方法，尤其后两种}。以下以柯布-道格拉斯生产函数为例
\begin{equation*}
    Y_t=K_t^\alpha\,L_t^{1-\alpha}
\end{equation*}
展示三种方法的原理和使用步骤。

\myskip
\noxmubluebf{方法一：直接方法}

对等式两侧去对数得
\begin{align*}
    \log Y_t & = \alpha\log K_t + \bka{1-\alpha}\log L_t\\ 
    \log \bar Y & = \alpha\log \bar K + \bka{1-\alpha}\log \bar L\label{eq:CDProdLog}
\end{align*}
两式相减
\begin{equation*}
    \log Y_t - \log \bar Y = \alpha\log K_t - \alpha\log \bar K + \bka{1-\alpha}\log L_t - \bka{1-\alpha}\log \bar L
\end{equation*}
定义对数偏离量，得
\begin{equation*}
    \hat y_t = \alpha\,\hat k_t + \bka{1-\alpha}\hat l_t\label{eq:CDProdLogLinear}
\end{equation*}

\myskip
\noxmubluebf{方法二：基础方法}

该方法要频繁使用到 \(X_t=\bar X\cdot e^{\hat x_t}\) 或近似结果 \(X_t=\bar X\bka{1+\hat x_t}\)，前者更常用。据此，有 \(X^\alpha_t=\bar X^\alpha\cdot e^{\alpha\,\hat x_t}\)（注意 \(\alpha\) 的位置，这也是常见性质），应用到生产函数中
\begin{align*}
    \bar Y\,e^{\hat y_t}=\bar K^\alpha\,e^{\alpha\,\hat k_t}\cdot \bar L^{1-\alpha}\,e^{\bka{1-\alpha}\hat l_t}
\end{align*}
此时可以取对数
\begin{align*}
    {\log\bar Y}+\hat y_t={\alpha\log\bar K}+\alpha\,\hat k_t+{\bka{1-\alpha}\log\bar L}+\bka{1-\alpha}\hat l_t
\end{align*}
消去稳态项即可得到结果；或使用 \(e^x\approx 1+x\) 近似
\begin{align*}
    \bar Y\bka{1+\hat y_t}&=\bar K^\alpha\,\bar L^{1-\alpha}\cdot\pa{1+\alpha\,\hat k_t}\ps{1+\pa{1-\alpha}\hat l_t}\\
    \bar Y\pa{1+\hat y_t}&=\bar K^\alpha\,\bar L^{1-\alpha}\cdot\ps{1+\alpha\,\hat k_t + \pa{1-\alpha}\hat l_t + \alpha\,\pa{1-\alpha}\underbrace{ \ \hat k_t\,\hat l_t \ }_{\approx \, 0}}
\end{align*}
出于对精度的容忍，丢弃高阶项也是使用本方法过程中的常用手段。剩下步骤省略。

\myskip
\noxmubluebf{方法三：通用方法}\label{sec:LogLinearGeneralMethod}

对数线性化的目标形如 \(Z_t=f\bka{X_t,Y_t}\)（或更多变量），首先对其取对数
\begin{align*}
    & \log Z_t=\log \, f\bka{e^{\log X_t}, \,e^{\log Y_t}}
\end{align*}
再求偏导数
\begin{align*}
    & \dfrac{\partial\log Z_t}{\partial\log X_t}\bka{\bar X,\bar Y}= \dfrac1{\bar Z}\cdot\dfrac{\p Z_t}{\p X_t}\bka{\bar X,\bar Y}\cdot \underbrace{ \ e^{\log \bar X} \ }_{=\,\bar X}=\underbrace{ \ \dfrac{\bar X}{\bar Z}\cdot\dfrac{\p Z_t}{\p X_t}\bka{\bar X,\bar Y} \ }_{\eqcolon\ \bar \eta_{zx}}
\end{align*}
此时再应用泰勒公式
\begin{align*}
    \log Z_t&\approx\log\bar Z+\dfrac{\bar X}{\bar Z}\dfrac{\p Z_t}{\p X_t}\bka{\bar X,\bar Y}\cdot\bka{\log X_t-\log\bar X}+\dfrac{\bar Y}{\bar Z}\dfrac{\p Z_t}{\p Y_t}\bka{\bar X,\bar Y}\cdot\bka{\log Y_t-\log\bar Y}
\end{align*}
最终得到两个形式
\begin{align*}
    \bar Z\,\hat z_t&=\dfrac{\p Z_t}{\p X_t}\bka{\bar X,\bar Y}\cdot\bar X\,\hat x_t+\dfrac{\p Z_t}{\p Y_t}\bka{\bar X,\bar Y}\cdot\bar Y\,\hat y_t\\
    \hat z_t&=\bar\eta_{zx}\cdot\hat x_t+\bar\eta_{zy}\cdot\hat y_t
\end{align*}
根据柯布-道格拉斯函数的经济含义，指数部分即为弹性，可立即得到结果。

使用通用方法需遵循：1. 计算导数并代入稳态；2. 写出泰勒展开式；3. 根据定义替换对数偏离量；4. 整理结果。

\subsection{对数线性化的性质}

下表总结了对数线性化的重要性质。
\begin{table}[H]
    \centering
    \caption{对数线性化的常用性质}
    \renewcommand{\arraystretch}{1.5}
    \begin{tabular}{>{\centering\arraybackslash}p{0.1\textwidth} >{\centering\arraybackslash}p{0.3\textwidth} >{\centering\arraybackslash}p{0.3\textwidth}}
        \toprule
        \textbf{序号} & \textbf{原方程} & \textbf{对数线性化形式} \\
        \midrule
        1 & \(X_t\,Y_t=c\) & \(\hat x_t+\hat y_t=0\) \\
        2 & \(X_t\islash Y_t=c\) & \(\hat x_t-\hat y_t=0\) \\
        3 & \(X_t\pm Y_t=c\) & \(\bar X\,\hat x_t\pm\bar Y\,\hat y_t=0\) \\
        \midrule
        4 & \(Z_t=c\,X_t\) & \(\hat z_t=\hat x_t\) \\
        5 & \(Z_t=X_t\,Y_t\) & \(\hat z_t=\hat x_t+\hat y_t\) \\
        6 & \(Z_t=X_t\islash Y_t\) & \(\hat z_t=\hat x_t-\hat y_t\) \\
        7 & \(Z_t=X_t^\alpha\) & \(\hat z_t=\alpha\,\hat x_t\) \\
        \midrule
        8 & \rule{0pt}{4ex}\(\dis Z_t=\sum\wsni\,X_{i,t}\) & \(\dis\hat z_t=\sum\wsni\,\frac{\bar X_{i}}{\bar Z}\cdot\hat x_{i,t}\) \\[2ex]
        \bottomrule
    \end{tabular}
\end{table}